We also ask whether two parameters per formality group can summarize the
eight size contrasts in the preferred categorical model.
\citet{FairDominguez1991} restrict coefficients on age-group population
shares to a quadratic polynomial and normalize them to sum to zero.
We adapt their polynomial restriction to firm-size coefficients, using log
assigned size rather than a category index. We retain solo work as the
reference instead of imposing a sum-zero normalization. This changes the
functional form, not the identification of the wage relationship.

Let \(z_i=\ln(\widetilde{Size}_i)\), using the assigned sizes in
Section~\ref{sec:methodology}. Write \(a_1\) and \(a_2\) for the informal
coefficients on log assigned size and its square. The corresponding
formal-minus-informal differences are \(b_1\) and \(b_2\), with \(k\)
indexing these two powers. The quadratic model is
\begin{equation}\label{eq:polynomial-size-profile}
  \ln(w_{iosdt})=\alpha+\phi F_i
  +\sum_{k=1}^{2}(a_k+b_k F_i)z_i^k
  +X_i'\gamma+\rho_o+\lambda_{s\times d\times t}
  +\varepsilon_{iosdt}.
\end{equation}
The sample, worker and occupational-position controls, fixed effects, weights,
and baseline clustering are the same as in Equation~(\ref{eq:main}). For status
\(f\in\{0,1\}\), write \(a_{kf}=a_k+b_k f\). The log-wage contrast
relative to solo work is \(q_f(z)=a_{1f}z+a_{2f}z^2\), which is zero at
the assigned solo size of one. Its implied elasticity is
\(a_{1f}+2a_{2f}z\), not a constant elasticity. The curve between assigned
sizes is imposed by the quadratic rather than observed in the data.

\begin{table}[htbp]
  \centering
  \caption{Quadratic size profiles with preferred controls}
  \label{tab:polynomial-size-profile}
  \small
  \begin{tabular}{lcc}
    \toprule
    Term & Formal workers & Informal workers \\[0.25em]
    \midrule
    Log assigned size ($\widehat{a}_{1f}$) & \shortstack{-0.0086\\{\scriptsize (0.0147)}} & \shortstack{0.1956***\\{\scriptsize (0.0134)}} \\[0.25em]
    Squared log assigned size ($\widehat{a}_{2f}$) & \shortstack{0.0089***\\{\scriptsize (0.0021)}} & \shortstack{-0.0251***\\{\scriptsize (0.0026)}} \\[0.25em]
    \midrule
    Implied 101+/solo premium (\%) & 19.8 & 41.9 \\[0.25em]
    \bottomrule
  \end{tabular}
  \vspace{0.3em}
  \begin{minipage}{0.95\textwidth}
  \footnotesize
  Notes: Both columns come from Equation~(\ref{eq:polynomial-size-profile}) on the preferred sample, with 101+ assigned 150 workers. The last row reports percentage premiums relative to solo work within each formality group. Parentheses contain sector-department-clustered standard errors. * $p<0.10$, ** $p<0.05$, *** $p<0.01$.
  \end{minipage}
\end{table}


Table~\ref{tab:polynomial-size-profile} reports the quadratic coefficients
when 101+ is assigned 150 workers. Figure~\ref{fig:polynomial-size-profile}
overlays the resulting curves and pointwise 95 percent confidence bands on
the preferred categorical estimates and their intervals. The informal curve
rises rapidly at small sizes and then bends downward, while the formal curve
has a much smaller initial decline followed by an increase. The quadratic
does not reproduce every categorical comparison. For formal 2--3 workers,
it predicts a gap of approximately zero relative to solo work, compared with -7.2 percent
in the categorical model. Its informal 101+/solo premium is 41.9 percent,
compared with 37.0 percent in the categorical model. The corresponding formal
101+/solo estimates are 19.8 and 17.2 percent.

\begin{figure}[htbp]
  \centering
  \includegraphics[width=0.95\textwidth]{preferred_quadratic_\figurelang.png}
  \caption{Quadratic and categorical firm-size profiles with preferred controls}
  \label{fig:polynomial-size-profile}
\end{figure}

Table~\ref{tab:quadratic-preferred-sensitivity} tests the functional-form
restriction and varies the size assigned to 101+. A quadratic constrains
eight within-formality categorical contrasts to depend on two coefficients,
imposing six restrictions per group and twelve jointly. We test those linear
restrictions using the unrestricted categorical estimates and their full
clustered covariance matrix. With 150 assigned to 101+, the sector-department
Wald \(F\)-statistics are 23.27 for formal workers and 22.34 for informal
workers, with six numerator and 558 denominator degrees of freedom.
Both tests reject the quadratic restriction at 1 percent, as does the joint
test.

\begin{table}[htbp]
  \centering
  \caption{Top-bin sensitivity and tests of the quadratic restriction}
  \label{tab:quadratic-preferred-sensitivity}
  \small
  \begin{tabular}{rrrrrr}
    \toprule
    Size for 101+ & Formal (\%) & Informal (\%) & $p_F$ & $p_I$ & $p_{Joint}$ \\[0.25em]
    \midrule
    \multicolumn{6}{l}{\textit{Sector-department clustering}} \\
    101 & 17.9 & 43.2 & $<0.001$ & $<0.001$ & $<0.001$ \\[0.25em]
    150 & 18.9 & 41.9 & $<0.001$ & $<0.001$ & $<0.001$ \\[0.25em]
    250 & 19.8 & 40.5 & $<0.001$ & $<0.001$ & $<0.001$ \\[0.25em]
    \midrule
    \multicolumn{6}{l}{\textit{Sector clustering}} \\
    101 & 17.9 & 43.2 & $<0.001$ & $<0.001$ & $<0.001$ \\[0.25em]
    150 & 18.9 & 41.9 & $<0.001$ & $<0.001$ & $<0.001$ \\[0.25em]
    250 & 19.8 & 40.5 & $<0.001$ & 0.001 & $<0.001$ \\[0.25em]
    \bottomrule
  \end{tabular}
  \vspace{0.3em}
  \begin{minipage}{0.95\textwidth}
  \footnotesize
  Notes: Formal and informal percentages are the implied 101+/solo premiums from the quadratic model, re-estimated for each assigned top-bin size. All other size assignments, observations, controls, and weights remain fixed. The $p$-values test whether the unrestricted categorical profiles from Equation~(\ref{eq:main}) lie on the specified quadratic. There are six restrictions per formality group and twelve jointly. Cluster-robust Wald statistics divided by the number of restrictions use an $F_{q,G-1}$ reference distribution. There are 559 sector-department clusters or 17 sector clusters. These are approximate tests, especially with only 17 clusters. The repeated percentages emphasize that changing clustering does not change point estimates. The value 150 is the baseline assignment, not an observed mean for 101+.
  \end{minipage}
\end{table}


Assigning 101, 150, or 250 workers to the open-ended category gives formal
101+/solo premiums of 18.7, 19.8, and 20.8 percent, respectively.
The informal counterparts are 43.2, 41.9, and 40.5 percent. The quadratic
restrictions are rejected at 1 percent under every assignment with either
sector-department or sector clustering. The latter tests use only 17 clusters
and should be read as approximate sensitivity checks. The shallow formal dip
under the baseline assignment lies between assigned sizes one and 2.5.
It disappears when 101+ is assigned 250, so it is not stable evidence of a
U-shaped conditional relationship. Confidence bands hold the size assignments
fixed and do not incorporate this uncertainty. We retain the categorical model as
the preferred specification.
\FloatBarrier
